\begin{abstract}
This paper studies nonlinear causal representation learning from an observational environment and one unknown-target perfect intervention for each scalar latent variable. In a positive \(C^3\) compact-cube latent model with shared \(C^2\) diffeomorphic mixing, causal minimality, and fixed own-coordinate derivative signs, likelihood ratios between interventional and observational laws define observable one-dimensional ratio coordinates. Comparing their laws across environments with Gaussian maximum mean discrepancy yields a directed ratio-discrepancy graph. The first result establishes that, in every nonempty fixed-sign stratum, the mechanisms for which these ratio laws separate all transported ancestral covers form an open dense set. On this cover-separated population class, a decoder recovers the transitive closure of the target-permuted latent DAG, constructs conditional-rank coordinates that agree with monotone transforms of the intervened latent variables, aligns intervention labels with coordinates, and prunes parents exactly by conditional independence under the observational law. Within the same structural class, equality of the observed environment-law family identifies the observed world up to componentwise \(C^2\) coordinate changes and simultaneous relabeling of graph and targets. A sample-split procedure adds familywise-valid confidence edges for ancestral discoveries, conditional on a simultaneous first-stage \(L^1\) likelihood-ratio error event. On the overlapping smooth compact-support faithful regime, the result gives an arbitrary-dimensional one-perfect-intervention-per-node counterpart to the bivariate law-separation route of \citet{vonKugelgenEtAl2023UnknownInterventions}.
\end{abstract}

\section{Introduction}

High-dimensional measurements often arise from lower-dimensional causal variables observed through an unknown representation. Causal representation learning asks when distributional variation across environments identifies those variables, their causal graph, and the intervention labels that generated the variation \citep{Scholkopf2021}. In nonlinear settings with invertible mixing, this question joins three traditions: graphical causal discovery from conditional independences \citep{Pearl2009,SpirtesGlymourScheines2000,Peters2017}, identifiable nonlinear independent component analysis from distributional changes \citep{Hyvarinen2016,Hyvarinen2019,Khemakhem2020,Hyvarinen2023}, and interventional discovery when targets are known, unknown, or partially specified \citep{Hauser2012,Yang2018,JaberEtAl2020SoftUnknownTargets,Dai2025,Zhang2026}.

The setting here is an observational law and \(n\) interventional laws generated by one perfect intervention for each of \(n\) scalar latent coordinates. The target labels are permuted, so the analyst observes environment laws before recovering the alignment between environments and latent variables. The latent variables live on the compact cube, the observation map is shared and diffeomorphic across environments, the local mechanisms are positive and smooth, and each intervention likelihood ratio has a fixed own-coordinate derivative sign. Under these conditions, the likelihood ratio \(R_i\), the density ratio for environment \(i\) against the observational law, is an observable scalar function whose latent form depends on the intervened coordinate and its parents.

The paper's central object is the law of one ratio coordinate under another environment. For each ordered pair \((j,i)\), the Gaussian maximum mean discrepancy \(D_{ji}\), the kernel discrepancy between the laws of \(R_i\) under environments \(0\) and \(j\), measures whether intervention \(j\) changes the distribution of ratio coordinate \(i\). \Cref{obj:thm:generic-cover-separation} shows that, within every nonempty positive smooth causal-minimal fixed-sign stratum, the class on which these discrepancies are positive for all transported ancestral covers is open and dense. The proof uses explicit regular witnesses from \cref{obj:prop:sparse-witness-certificate} and analytic perturbations along normalized affine paths to move edge-specific second-moment contrasts away from cancellation.

On that generic population class, \cref{obj:thm:exact-ratio-decoder} gives an exact decoder. Starting from the environment laws, \(\mathscr D\), the population decoder, forms the ratio-discrepancy graph, chooses a topological order, constructs conditional distribution transforms of log-ratios, aligns each environment with the coordinate recovered by its own rank, and prunes candidate predecessors using conditional independence under the observational law. The ratio graph recovers the transitive closure of \(G^\pi\), the target-permuted DAG; the conditional ranks equal \(Q_{\pi(i)}(V_{\pi(i)})\), the intervention distribution transform of the target latent coordinate, or its sign-reversed version; and causal minimality makes the true parent set the unique inclusion-minimal admissible set. Thus the output graph is exactly \(G^\pi\), the latent DAG expressed on environment labels.

The result contributes to the unknown-intervention literature by using likelihood-ratio laws in place of scores, support geometry, quadratic precision structure, or supplied target metadata. The closest invertible-mixing comparator is \citet{vonKugelgenEtAl2023UnknownInterventions}, who establish a bivariate one-intervention result under a continuous law-separation condition and an arbitrary-dimensional result with two paired perfect interventions per node. On the overlapping positive smooth compact-support faithful regime, \cref{obj:thm:exact-ratio-decoder} supplies an arbitrary-dimensional one-perfect-intervention-per-node population recovery statement and gives the Gaussian-MMD law-separation witness in the two-variable case. Relative to target-aligned approaches such as \citet{WendongEtAl2023CauCA} and \citet{YaoEtAl2025InvariancePrinciple}, the ratio construction recovers the graph, order, and target alignment from the same observed law family.

The statistical layer translates population discrepancies into finite-sample edge statements. \Cref{obj:thm:simultaneous-confidence-edges} studies a sample-split design: independent training folds produce fitted positive likelihood ratios, and independent evaluation folds compute empirical Gaussian-MMD discrepancies. Given a simultaneous first-stage \(L^1\) ratio-error radius \(a_N(\eta)\), the theorem provides a familywise event on which every selected edge in the confidence graph has an ancestral interpretation in the latent DAG. When every transported ancestral-cover discrepancy exceeds twice the simultaneous radius, the selected graph has the same transitive closure as \(G^\pi\). This connects the ratio decoder to kernel two-sample concentration \citep{GrettonEtAl2012KernelTwoSample,SchrabEtAl2023MMDAgg} while keeping the first-stage ratio-estimation requirement explicit.

The final section formulates the next coordinate-error problem in a bounded Hölder regime. Under smoothness, lower-envelope, own-derivative margin, shared-mixing geometry, observed log-ratio regularity, and conditional-design assumptions, \cref{obj:oeq:generated-rank-frontier} states a generated-rank rate target for conditional-CDF estimation with fitted log-ratio responses and fitted conditioning covariates. The target combines second-stage smoothing bias, local empirical variation, and first-stage \(C^1\) error propagated through bandwidth neighborhoods, drawing on conditional-CDF and generated-regressor theory \citep{Xie2023UniformConditionalCDF,MammenRotheSchienle2012GeneratedCovariates,Lee2018PartialMeanGeneratedRegressors}.

The whole argument rests on two observations that the rest of the paper makes exact. Under the stated positivity, smoothness, causal-minimality, shared-diffeomorphic-mixing, one-perfect-intervention, fixed-sign, and cover-separation conditions, the ratios \(R_i\), the density ratios for interventional environments against the observational law, support ancestry recovery through population discrepancies, and the conditional-rank construction returns a monotone transform of the corresponding latent coordinate. Together with causal minimality, shared diffeomorphic mixing, and the one-perfect-intervention-per-node setup, positivity, smoothness, the sign condition, and cover separation make those two steps exact; the bookkeeping objects introduced below state these conditions precisely. A reader who wants the argument before its apparatus can read \cref{obj:def:population-decoder} and \cref{obj:thm:exact-ratio-decoder} first and consult \cref{sec:setup} for the objects they use. The formal layer accompanying the paper is machine-checked, with precise scope recorded in \cref{sec:appendix}.

\Cref{sec:related-work} situates the paper in graphical causal inference, nonlinear ICA, causal representation learning, unknown-target interventions, and kernel testing. \Cref{sec:setup} introduces the positive compact-cube model, observed worlds, likelihood ratios, and ratio-law discrepancies. \Cref{sec:generic-separation} proves generic ratio separation. \Cref{sec:population-decoding} gives the population decoder and its exact recovery theorem. \Cref{sec:confidence-edges} establishes sample-split confidence edges, and \cref{sec:limitations} records the generated-rank frontier and further directions.
